%----------------------------------------------------------------------
\section{Category Theory from Discrimination}\label{sec:cattheory}
%----------------------------------------------------------------------

Section \ref{sec:homotopy} constructed a symmetric $n$-groupoid\footnote{Section-level Toulmin. \emph{Move}: add direction via $\tau/\sigma$ asymmetry. \emph{ZFC comparison}: category theory is typically layered on top of set theory as a separate framework. Here, categories are constructed from the same algebra, with direction from datum properties. The native structure is 2-categorical; the 1-categorical shadow is a principled truncation whose residue (the lost composition data) is tracked, not discarded. \emph{Residue instance}: the 2-cell $d_1 \wedge d_2$ records which discriminations were composed. Projecting to the 1-categorical shadow loses this; the loss is itself a residue.}
from the exterior algebra over $\GF(2)$
(plus profile linearity (Definition \ref{def:profile-linear})). In a groupoid, every
morphism is invertible. General category theory requires
\emph{directed} morphisms that need not be invertible. We now
show that direction arises from the $\tau/\sigma$ asymmetry
of the discrimination process, layered on top of the symmetric
algebra.

\subsection{Direction from $\tau/\sigma$ asymmetry}

\begin{remark}[The source of asymmetry]
The exterior algebra over $\GF(2)$ is commutative:
$d_1 \wedge d_2 = d_2 \wedge d_1$. This symmetry produces
groupoids. But the discrimination \emph{process} is not
symmetric: a discriminator $d$ assigns $\tau_d(a) = 1$,
$\sigma_d(b) = 1$ for elements $a, b$ in a population.
The labels $\tau$ and $\sigma$ are distinct; swapping them
produces a different discrimination. This labeling asymmetry,
absent from the algebra, is what introduces direction.
\end{remark}

\begin{definition}[Directed morphism]\label{def:directed-morphism}
Let $S$ be a population under regime $\kappa$, and let
$a, b \in S$. A \emph{directed morphism} $f: a \to b$ is a
discriminator $d \in \kappa$ such that $\tau_d(a) = 1$ and
$\sigma_d(b) = 1$. The element $a$ is the \emph{source}
(classified under $\tau$) and $b$ is the \emph{target}
(classified under $\sigma$).

Two directed morphisms $f: a \to b$ and $g: a \to b$
witnessed by distinct discriminators $d_f \neq d_g$ are
\emph{distinct morphisms} between the same source and target.
\end{definition}

\begin{proposition}[Non-invertibility]\label{prop:noninvert}
A directed morphism $f: a \to b$ witnessed by discriminator
$d$ does not automatically yield a morphism $b \to a$.
The reverse morphism requires a \emph{separate} discriminator
$d'$ with $\tau_{d'}(b) = 1$ and $\sigma_{d'}(a) = 1$.
This $d'$ is an independent dimension of $\kappa$.
\end{proposition}

\begin{proof}
The discriminator $d$ assigns specific $\tau/\sigma$ values
to elements of $S$. These assignments are fixed properties
of $d$. Over $\GF(2)$, $d = -d$, so the algebraic inverse
of $d$ is $d$ itself --- but algebraic self-inverse does not
swap the $\tau$ and $\sigma$ labels. The labels are structural
metadata of the discrimination, not elements of $\GF(2)$.
Reversing the direction $a \to b$ to $b \to a$ requires
a discriminator with the opposite labeling, which is a
distinct element of $\kappa$.
\end{proof}

\begin{corollary}[Cost of isomorphism]\label{cor:iso-cost}
An isomorphism between $a$ and $b$ requires two discriminators:
one witnessing $a \to b$ and one witnessing $b \to a$.
Isomorphisms cost two $\kappa$-dimensions; mere morphisms
cost one. Invertibility is not free; it is purchased by
additional discrimination.
\end{corollary}

\subsection{The constructed category $\mathcal{C}(S,\kappa)$}

\begin{definition}[Emergent category]\label{def:category}
Let $S$ be a population and $\kappa$ a discrimination regime.
The \emph{category} $\mathcal{C}(S,\kappa)$ has:
\begin{itemize}[nosep]
  \item \textbf{Objects}: $\kappa$-equivalence classes of
    elements of $S$.
  \item \textbf{Morphisms}: directed morphisms
    (Definition \ref{def:directed-morphism}). For objects
    $A, B$ (equivalence classes), $\operatorname{Hom}(A,B)$
    is the set of discriminators $d \in \kappa$ such that
    $\tau_d(a) = 1$ for all $a \in A$ and $\sigma_d(b) = 1$
    for all $b \in B$.
  \item \textbf{Identity}: for each object $A$, the identity
    morphism $\mathrm{id}_A$ is the degenerate discriminator
    $d \wedge d = 0$ (Definition \ref{def:degeneracy}).
    Self-discrimination produces no partition: $A$ is mapped
    to itself trivially.
  \item \textbf{Composition}: defined below.
\end{itemize}
\end{definition}

\begin{definition}[Composition via filtration]\label{def:composition}
Let $f: A \to B$ be witnessed by $d_1$ and $g: B \to C$ be
witnessed by $d_2$. The composite $g \circ f$ is
the joint filtration of $S$ through $d_1$ and $d_2$,
recorded as the grade-2 element
$d_1 \wedge d_2 \in \bigwedge^2 V$.

This is a \emph{2-cell}, not a 1-cell. The composite of
two directed 1-morphisms produces a 2-morphism witnessing
their joint action. The $\tau/\sigma$ assignments track
source ($A \in \tau_{d_1}$) and target
($C \in \sigma_{d_2}$) through the filtration chain, with
$B$ as the intermediate object ($\sigma_{d_1}(B) = 1$,
$\tau_{d_2}(B) = 1$).
\end{definition}

\begin{remark}[The interchange law and self-hosting]
\label{rem:interchange-selfhosting}
The interchange law is not merely a coherence condition
for the 2-category. It is the structural backbone of
self-hosting (Theorem \ref{thm:self-hosting}) and the
halting tower (Remark \ref{rem:hyper-recursive}).

Horizontal composition ($*$) is sequential
$\kappa$-evolution: applying discriminators in sequence.
Vertical composition ($\circ$) is perspective shift:
the $S_3$ rotation that reinterprets which component
plays which role. The interchange law states that
these commute:
\[
  (\beta \circ \alpha) * (\delta \circ \gamma)
  = (\beta * \delta) \circ (\alpha * \gamma).
\]
In self-hosting terms: the outer system models an inner
system. From outside, the inner system's behavior is
finite and comprehensible (a $\kappa$-regime we
observe). From inside, the outer world is unfathomable
(the inner system's halting problem). The self-hosting
isomorphism (inner $\cong$ outer) transfers results
between levels. The interchange law says this transfer
commutes with $\kappa$-evolution: what you learn by
modeling the inner system's epistemic gap can be
transferred to characterize the outer system's own gap,
because evolving $\kappa$ and shifting perspective
produce the same result in either order.

Without interchange, the level-invariance of the halting
tower (Remark \ref{rem:hyper-recursive}) would be a
claim about each level separately. With interchange, it
is a structural fact: the gap between ``finite from
outside'' and ``unfathomable from inside'' has the
same shape at every level because the two operations
that define ``level'' (evolution and perspective shift)
commute. The explicit verification (Appendix B.20)
shows both paths through the interchange square yield
$p(g{\circ}f) + p(g'{\circ}f') \in \GF(2)^2$: the
cross terms cancel because $1+1=0$, collapsing the
grade-4 expression to grade 2 via the Klein
four-group's characteristic 2. The same $1+1=0$ that
excludes self-reference, kills Russell's affine
offset, and makes the halting tower level-invariant
also makes the interchange law hold.
\end{remark}

\begin{remark}[Interchange and Russell are the same claim]
\label{rem:interchange-russell}
The interchange law (Theorem \ref{thm:2cat}) and
Russell exclusion (Theorem \ref{thm:russell}) share
the technique profile $[1{+}1{=}0 + \text{Linear}]$
(Appendix \ref{sec:unified}, Claim F). Russell:
evaluation linearity rejects the affine predicate
($\tau(0) = 1 \neq 0$). Interchange: evaluation
linearity kills the quadratic cross term ($2 = 0$).
Both are the same fact: linearity over $\GF(2)$
rejects the non-linear component. Self-reference
and incoherence are the same pathology, excluded
by the same axiom.
\end{remark}

\begin{remark}[The framework is 2-categorical from the start]
\label{rem:2cat}
Composition of 1-morphisms producing a 2-cell means
$\mathcal{C}(S,\kappa)$ is not a 1-category but a
\emph{2-category from the start}: objects are
$\kappa$-equivalence classes, 1-morphisms are
discriminators (grade-1 elements), and composition of
1-morphisms is a 2-morphism (grade-2 element). This is
not a deficiency but a feature: the framework entangles
direction and higher structure from the beginning,
rather than constructing a 1-category and then adding
higher cells. The exterior algebra's grading \emph{is}
the categorical dimension.

A 1-categorical shadow is obtained by projecting: identify
two 1-morphisms $f, g: A \to B$ whenever they are
connected by a 2-cell (i.e.\ $d_f \wedge d_g \neq 0$).
The resulting quotient is a 1-category, but the 2-cells
carry information that the quotient loses.
(Full 2-cell analysis: \S\ref{app:composition-2cell}.)
\end{remark}

\begin{theorem}[2-category axioms]\label{thm:cat-axioms}\footnote{Axiom class \textsf{A} $\cdot$ depth 4 $\cdot$ \S\ref{app:2cat-full}.  \emph{Warrant}: associativity and commutativity of $\wedge$; interchange law holds. \emph{Qualifier}: the native structure is 2-categorical; the 1-categorical shadow is a quotient Finite-$\kappa$: yes. Novel (composition produces 2-cells natively). (\S\ref{app:composition-2cell}).}
$\mathcal{C}(S,\kappa)$ satisfies the axioms of a
strict 2-category.
\end{theorem}

\begin{proof}
\emph{Horizontal composition (associativity).}
Let $f: A \to B$, $g: B \to C$, $h: C \to D$ be witnessed
by $d_1, d_2, d_3$ respectively. Then
$(h \circ g) \circ f$ is witnessed by
$d_1 \wedge (d_2 \wedge d_3)$ and
$h \circ (g \circ f)$ is witnessed by
$(d_1 \wedge d_2) \wedge d_3$. By associativity of the
exterior algebra, these are equal. The $\tau/\sigma$
assignments are tracked through the filtration chain,
and the monoidal property (Proposition \ref{prop:monoidal})
guarantees that the order of filtration does not affect
the result.

\emph{Identity.}
For any $f: A \to B$ witnessed by $d$, the composite
$f \circ \mathrm{id}_A$ is witnessed by $d \wedge (d_A \wedge d_A)
= d \wedge 0 = d$ (where $d_A$ is any discriminator
associated with $A$). Similarly $\mathrm{id}_B \circ f = d$.
The identity morphism, being the degenerate self-discrimination
(grade 0), acts as the unit for the wedge product.
The interchange law holds by associativity and
commutativity of $\wedge$ over $\GF(2)$.
(Full axiom check including interchange: \S\ref{app:2cat-full}.)
\end{proof}

\subsection{Functors as $\kappa$-deformations}\label{sec:functors}

\begin{definition}[Functor]\label{def:functor}
Let $\kappa$ and $\kappa'$ be two discrimination regimes
(possibly over different populations $S$ and $S'$). A
\emph{functor} $F: \mathcal{C}(S,\kappa) \to
\mathcal{C}(S',\kappa')$ is a $\kappa$-deformation: a
structure-preserving map that sends:
\begin{itemize}[nosep]
  \item discriminators in $\kappa$ to discriminators in
    $\kappa'$,
  \item $\tau/\sigma$ assignments under $\kappa$ to
    $\tau/\sigma$ assignments under $\kappa'$,
\end{itemize}
such that the wedge product and boundary structure are
preserved.
\end{definition}

\begin{proposition}[Forgetful and free functors]\label{prop:forgetful}
Let $\kappa' \subset \kappa$ be a sub-regime.
\begin{enumerate}[label=(\roman*),nosep]
  \item The projection $\pi: \kappa \to \kappa'$ (omitting
    dimensions) induces a \emph{forgetful functor}
    $U: \mathcal{C}(S,\kappa) \to \mathcal{C}(S,\kappa')$.
    $U$ sends each object to its coarser equivalence class
    (merging objects that were distinguished only by the
    omitted dimensions) and each morphism to its image under
    projection (forgetting the contribution of omitted
    discriminators).
  \item The inclusion $\iota: \kappa' \hookrightarrow \kappa$
    induces a \emph{free functor}
    $F: \mathcal{C}(S,\kappa') \to \mathcal{C}(S,\kappa)$.
    $F$ sends each object to itself (now embedded in a
    finer equivalence structure) and each morphism to itself
    (now admitting additional morphisms that the finer
    $\kappa$ makes available).
\end{enumerate}
\end{proposition}

\begin{proof}
We verify the functor axioms explicitly for the forgetful
functor $U$ (the free functor $F$ is analogous).

\emph{Composition preservation.}
Let $f: A \to B$ and $g: B \to C$ be morphisms in
$\mathcal{C}(S,\kappa)$, witnessed by discriminators
$d_f, d_g \in \kappa$. Their composite $g \circ f$ is
the joint filtration, whose algebraic content is recorded
as $d_f \wedge d_g$ with $\tau/\sigma$ tracking
(Definition \ref{def:composition}).

The forgetful functor $U$ applies the projection
$\pi: \bigwedge(\GF(2)^n) \to \bigwedge(\GF(2)^{n'})$
(where $n' < n$ corresponds to the smaller regime
$\kappa'$). This projection is a \emph{ring homomorphism}
of exterior algebras: it preserves both addition and the
wedge product. Therefore:
\[
  U(g \circ f) = \pi(d_f \wedge d_g)
  = \pi(d_f) \wedge \pi(d_g)
  = U(f) \wedge U(g) = U(g) \circ U(f).
\]
The second equality is the ring homomorphism property.
The $\tau/\sigma$ tracking is preserved because $\pi$
acts on each discriminator independently, and the
$\tau/\sigma$ assignments of $\pi(d)$ are the restrictions
of $d$'s assignments to the dimensions retained in
$\kappa'$.

If a discriminator $d$ involves only dimensions outside
$\kappa'$, then $\pi(d) = 0$ (the discriminator is
projected to the identity). This is correct: the morphism
it witnessed is invisible at the coarser resolution,
and $U$ maps it to the identity morphism.

\emph{Identity preservation.}
The identity morphism $\mathrm{id}_A$ is the degenerate
discriminator $d \wedge d = 0$.
$U(\mathrm{id}_A) = \pi(0) = 0 = \mathrm{id}_{U(A)}$.
\end{proof}

\begin{theorem}[Triadic adjunction system]\label{thm:adjunction}\footnote{Axiom class \textsf{A} $\cdot$ depth 5 $\cdot$ \S\ref{app:adjunction-full}.  \emph{Warrant}: $\pi \circ \iota = \mathrm{{id}}$ (direct summand). \emph{Qualifier}: these are retract adjunctions --- individually trivial, collectively non-trivial as covering families for the Grothendieck topology. Finite-$\kappa$: yes. Standard (retract adjunction from direct summand). Collectively novel (six adjunctions as covering families).}
The $S_3$ symmetry of the triangle identity generates six
adjunctions --- three chirality pairs --- by permuting
$\kappa$, $\sigma$, $\tau$ over the adjunction stencil
$A \leftarrow B \to C$, where $A$ is the left adjoint,
$B$ is the hom, and $C$ is the right adjoint.
\end{theorem}

\begin{proof}
The adjunction stencil $A \leftarrow B \to C$ places three
roles in a fixed relationship: the hom-object $B$ mediates
between left adjoint $A$ and right adjoint $C$. The
adjunction is the natural bijection
$\operatorname{Hom}(A, C) \cong B$, expressing that $B$
determines the morphism space between $A$ and $C$.

The triangle identity (Definition \ref{def:triangle})
states that any two of $\{\kappa, \sigma, \tau\}$ determine
the third. This is precisely the structure of an adjunction:
given the left and right adjoints, the hom is determined, and
given the hom and one adjoint, the other is determined.

Permuting $\kappa, \sigma, \tau$ over the stencil
$A \leftarrow B \to C$ yields six assignments, grouped
into three pairs of opposite chirality:

\emph{On the $\sigma$-axis} ($\sigma$ as hom):
\begin{itemize}[nosep]
  \item $\kappa \leftarrow \sigma \to \tau$:
    $\kappa$ is left adjoint, $\sigma$ is the hom-space,
    $\tau$ is right adjoint.
  \item $\tau \leftarrow \sigma \to \kappa$:
    chirality reversed. $\tau$ is left adjoint,
    $\kappa$ is right adjoint.
\end{itemize}

\emph{On the $\tau$-axis} ($\tau$ as hom):
\begin{itemize}[nosep]
  \item $\kappa \leftarrow \tau \to \sigma$:
    $\kappa$ is left adjoint, $\tau$ is the hom-space,
    $\sigma$ is right adjoint.
  \item $\sigma \leftarrow \tau \to \kappa$:
    chirality reversed.
\end{itemize}

\emph{On the $\kappa$-axis} ($\kappa$ as hom):
\begin{itemize}[nosep]
  \item $\sigma \leftarrow \kappa \to \tau$:
    $\sigma$ is left adjoint, $\kappa$ is the hom-space,
    $\tau$ is right adjoint.
  \item $\tau \leftarrow \kappa \to \sigma$:
    chirality reversed.
\end{itemize}

Each assignment produces a valid adjunction because the
triangle identity guarantees the required bijection. We
prove the bijection for the first assignment; the others
follow by $S_3$ symmetry.

\emph{Bijection.}
For the $\sigma$-axis ($\kappa \leftarrow \sigma \to \tau$):
a morphism from a $\kappa$-element to a $\tau$-element in
the $\sigma$-regime is determined by its $\GF(2)$-value
via $\sigma = \kappa + \tau$. The maps
$\Phi: \operatorname{Hom}_\sigma(\kappa, \tau) \to \sigma$
(extract the $\sigma$-value) and
$\Psi: \sigma \to \operatorname{Hom}_\sigma(\kappa, \tau)$
(reconstruct from the value via $\kappa + s = \tau$) are
mutual inverses because $\GF(2)$ addition is bijective.
The bijection is:
\begin{equation}\label{eq:triadic-adj}
  \operatorname{Hom}_\sigma(\kappa, \tau) \cong \sigma.
\end{equation}

\emph{Naturality.}
$\Phi$ is natural in both variables because
$\pi_\sigma$ (projection onto $\sigma$-components) is a
ring homomorphism preserving $\wedge$ and $+$.
Precomposition with $F(f) = \pi_\sigma(d_f)$ and
post-composition with $U(g) = \pi_\tau(d_g)$ both
commute with $\Phi$ because $\Phi$ extracts the
$\sigma$-value, and ring homomorphisms preserve
products.

\emph{Zig-zag equations.}
$\kappa'$ is a direct summand of $\kappa$ as a
$\GF(2)$-vector space ($V = \kappa' \oplus W$), so
$\pi \circ \iota = \mathrm{id}_{\kappa'}$. The unit
$\eta = \mathrm{id}$ and the counit $\varepsilon =
\mathrm{id}$ on the image. Both zig-zag equations
reduce to $\mathrm{id} \circ \mathrm{id} = \mathrm{id}$.
This is a retract adjunction: structurally simple but
algebraically genuine.

\emph{Six adjunctions.}
The remaining five follow by $S_3$ permutation of
$\{\kappa, \sigma, \tau\}$ over the stencil. Within
each axis, the two chiralities ($\kappa \leftarrow
\sigma \to \tau$ and $\tau \leftarrow \sigma \to \kappa$)
are related by reversing left and right adjoints. The
triangle identity ensures all six access the same
information ($\dim(\kappa) = \dim(\sigma) = \dim(\tau)$,
preserved by XOR). (Full component verification: \S\ref{app:adjunction-full}.)
\end{proof}

\begin{corollary}[The free-forgetful adjunction is one of six]
\label{cor:free-fgt}
The free-forgetful adjunction between coarsening
($\pi: \kappa \to \kappa'$) and refinement
($\iota: \kappa' \hookrightarrow \kappa$) is the
$\kappa$-axis adjunction $\sigma \leftarrow \kappa \to \tau$:
$\kappa$ plays the hom role (the morphism space is the
set of discriminators), $\sigma$ is the free (left) functor
(embedding into a finer regime), and $\tau$ is the
forgetful (right) functor (projecting to a coarser regime).

The remaining five adjunctions are equally valid:
\begin{itemize}[nosep]
  \item The $\sigma$-axis adjunctions govern the relationship
    between population-level structure ($\kappa$) and
    observation-level structure ($\tau$), mediated by
    meta-discrimination ($\sigma$).
  \item The $\tau$-axis adjunctions govern the relationship
    between population-level structure ($\kappa$) and
    meta-discrimination ($\sigma$), mediated by
    observation ($\tau$).
\end{itemize}
\end{corollary}

\begin{remark}
The six adjunctions are not six separate constructions.
They are one structure --- the $\GF(2)$ toroid --- viewed
from six orientations. Each orientation places one
component in the hom role and the other two in adjoint
roles, with chirality determining which is left and which
is right. The adjunction is the toroid.
\end{remark}

\subsection{Natural transformations}\label{sec:nattrans}

\begin{definition}[Natural transformation]\label{def:nattrans}
Let $F, G: \mathcal{C}(S,\kappa) \to \mathcal{C}(S',\kappa')$
be two functors (two $\kappa$-deformations between the same
source and target categories). A \emph{natural transformation}
$\eta: F \Rightarrow G$ is a family of morphisms
$\{\eta_A: F(A) \to G(A)\}_{A \in \mathrm{Ob}(\mathcal{C}(S,\kappa))}$
in $\mathcal{C}(S',\kappa')$ such that for every morphism
$f: A \to B$ in $\mathcal{C}(S,\kappa)$, the following
square commutes:
$G(f) \circ \eta_A = \eta_B \circ F(f)$.

In the discrimination framework, each component $\eta_A$ is
a discriminator in $\kappa'$ that mediates between the
two deformations' images of $A$. The naturality condition
states that mediation commutes with the deformations'
action on morphisms.
\end{definition}

\begin{proposition}[Naturality from coherence of discrimination]
\label{prop:naturality}
Let $F$ and $G$ be two $\kappa$-deformations that differ on
a set of dimensions $D \subset \kappa$. A natural
transformation $\eta: F \Rightarrow G$ exists if and only if
there is a coherent family of discriminators in $\kappa'$
--- one per object of the source category --- that
compensates for the difference between $F$ and $G$ on the
dimensions in $D$, and this compensation commutes with
all morphisms in the source category.
\end{proposition}

\begin{proof}
The two deformations $F$ and $G$ agree on dimensions
outside $D$ and disagree on dimensions in $D$. For each
object $A$, $F(A)$ and $G(A)$ differ in their $\tau/\sigma$
profiles on the dimensions in $D$. A discriminator
$\eta_A \in \kappa'$ with $\tau_{\eta_A}(F(A)) = 1$ and
$\sigma_{\eta_A}(G(A)) = 1$ provides the component morphism.

The coherence condition (naturality) requires that for
$f: A \to B$ witnessed by $d$, the composite discriminations
$G(d) \wedge \eta_A$ and $\eta_B \wedge F(d)$ agree. Since
the wedge product is commutative over $\GF(2)$, this reduces
to verifying that the $\tau/\sigma$ tracking is consistent
through both paths around the square. This holds when the
family $\{\eta_A\}$ is determined by the dimensions $D$
uniformly --- that is, when the same dimensional shift
applies at every object.
\end{proof}

\subsection{The Yoneda lemma}\label{sec:yoneda}

The Yoneda lemma is the statement that an object is fully
determined by its morphism profile. In the discrimination
framework, this is not a lemma to be proved but a
\emph{design principle that has been built in from the
start}. We now show that the formal Yoneda lemma is a
theorem of the framework.

\begin{definition}[Representable presheaf]\label{def:representable}
For an object $A$ in $\mathcal{C}(S,\kappa)$, the
\emph{representable presheaf} $h_A = \operatorname{Hom}(-, A)$
assigns to each object $B$ the set of discriminators
$d \in \kappa$ with $\tau_d(B) = 1$ and $\sigma_d(A) = 1$.
This is the set of all directed morphisms into $A$.
\end{definition}

\begin{theorem}[Yoneda]\label{thm:yoneda}\footnote{Axiom class \textsf{A} $\cdot$ depth 5 $\cdot$ \S\ref{app:yoneda-full}.  \emph{Warrant}: standard for any locally small category. \emph{Qualifier}: local smallness is immediate for finite $\kappa$; for infinite structures, it inherits the operationalist axiom's strength or weakness Finite-$\kappa$: yes. Re-description. Local smallness caveat for infinite. (\S6, local smallness caveat).}
The Yoneda embedding
$\mathcal{Y}: \mathcal{C}(S,\kappa) \hookrightarrow
[\mathcal{C}(S,\kappa)^{\mathrm{op}}, \mathbf{Set}]$
defined by $A \mapsto h_A$ is full and faithful. That is:
$A \cong B$ in $\mathcal{C}(S,\kappa)$ if and only if
$h_A \cong h_B$ as presheaves.
\end{theorem}

\begin{proof}
\emph{Faithfulness.}
If $f \neq g: A \to B$ (distinct discriminators $d_f \neq d_g$),
then $h_f \neq h_g$ because $d_f$ and $d_g$ produce
different $\tau/\sigma$ assignments on at least one element
of $S$ (they are distinct dimensions of $\kappa$), so
there exists an object $C$ for which
$h_f(C) \neq h_g(C)$.

\emph{Fullness.}
Let $\eta: h_A \to h_B$ be a natural transformation. By
the Yoneda argument, $\eta$ is determined by its component
at $A$: $\eta_A \in h_B(A) = \operatorname{Hom}(A, B)$.
This component is a discriminator $d$ with $\tau_d(A) = 1$
and $\sigma_d(B) = 1$. The morphism $d: A \to B$ in
$\mathcal{C}(S,\kappa)$ maps to $\eta$ under $\mathcal{Y}$.

\emph{Equivalently, in the language of the framework:}
two objects are isomorphic if and only if they have
identical morphism profiles against the entire population
--- that is, no discriminator in $\kappa$ separates them
directionally. This is $\kappa$-extensionality
(Theorem \ref{thm:ext}) restricted to directed morphisms.
The Yoneda embedding is the directed version of the
principle that identity is constituted by discrimination
profile.
(Full proof: \S\ref{app:yoneda-full}.)
\end{proof}

\begin{remark}
That Yoneda is a theorem rather than a surprise is the point.
The framework was constructed so that an element's identity is
its relational profile (Remark after Definition
\ref{def:membership}). Yoneda formalizes this in the directed
setting and confirms that the construction is consistent:
the categorical structure constructed from discrimination
inherits the relational-identity principle from which it
was built.

For finite $S$ and $\kappa$, local smallness (that
$\operatorname{Hom}(A,B)$ is a set) is immediate. For
infinite structures obtained via the operationalist
axiom, the set-theoretic status of hom-sets is exactly
the kind of question the framework is addressing, and
there is a circularity risk: Yoneda presupposes that
hom-sets are well-behaved, but their well-behavedness
depends on the framework that Yoneda helps to validate.
The finite case is clean; the infinite case inherits
whatever strength or weakness Finite-$\kappa$: yes. Re-description. Local smallness caveat for infinite. the operationalist axiom (Definition \ref{def:operationalist})
confers.
\end{remark}

\subsection{Limits and colimits from residue structure}

\begin{definition}[Limit as universal cone]\label{def:limit}
Let $D: \mathcal{J} \to \mathcal{C}(S,\kappa)$ be a diagram
(a functor from a small index category). A \emph{limit} of
$D$ is an object $L$ with morphisms $\pi_j: L \to D(j)$
for each $j \in \mathcal{J}$, universal among all such cones.

In the discrimination framework, $L$ is the $\kappa$-equivalence
class that is maximally classified under $\tau$ by every
discriminator witnessing a morphism into any $D(j)$. The
universality condition states: $L$ is the tightest equivalence
class from which every object in the diagram is reachable.
\end{definition}

\begin{proposition}[Limits from exhaustive filtration]
\label{prop:limits}
The limit of a diagram $D$ exists in $\mathcal{C}(S,\kappa)$
whenever $\kappa$ contains sufficient discriminators to
express the universal property. If the limit does not exist,
the obstruction is a residue: the wedge product of the cone
morphisms produces a non-zero element in the kernel, and
consuming this residue via $\kappa$-evolution extends the
category to one in which the limit exists.
\end{proposition}

\begin{proof}
The cone condition requires, for each pair of objects
$D(j), D(k)$ connected by a morphism $D(f): D(j) \to D(k)$,
that $\pi_k = D(f) \circ \pi_j$. This is a system of
equations in the discriminators of $\kappa$, mediated by
the wedge product. If the system has a solution --- a
discriminator (or grade-$k$ element) that simultaneously
satisfies all cone conditions --- the limit exists.

If no solution exists, the wedge product of the cone
constraints produces a non-trivial residue: the
constraints are mutually incompatible under the current
$\kappa$. By Theorem \ref{thm:exhaustive} (operational
exhaustivity), this residue specifies the discriminator
needed to resolve the incompatibility. Consuming the
residue evolves $\kappa$, and in the evolved regime
the limit exists.
\end{proof}

\begin{remark}
The dual construction gives colimits: the colimit of a
diagram is the loosest equivalence class into which every
object of the diagram maps. Colimits that fail to exist
produce residues with the same self-specifying property.
\end{remark}

\begin{remark}[Completeness is not a fixed property]
A category $\mathcal{C}(S,\kappa)$ is complete (has all
small limits) relative to its current $\kappa$-state.
Introducing a new diagram that lacks a limit triggers
$\kappa$-evolution, extending the category to accommodate
it. Completeness is not a static property of the category
but a dynamic relationship between the category and the
demands placed on it. Every demand that exposes an
incompleteness simultaneously provides (via the residue)
the means to resolve it.
\end{remark}

